% \aj{Perhaps we can change the title to something like "Background and Model" as we are also introducing how we model the loss etc.}

% \aj{One important thing missing that's missing is the non-parametric representation of the distributions. Right now, it is pushed into the proof of theorem $1$ in the supplementary material \cref{sec:sard-pol-grad}}
\subsection{Reinforcement Learning from Human Feedback (RLHF)}

The most popular method of aligning LLMs with human preferences is Reinforcement Learning from Human Feedback (RLHF). The standard RLHF pipeline consists of three phases as described in \cite{Christiano2017DeepRL,Ouyang2022TrainingLM}. The first phase is a supervised fine-tuning (SFT) phase where the outputs of a pre-trained model are aligned with human responses. This is followed by a reward modeling phase, where a reward model is trained using human preferences and finally, a reinforcement learning phase where the model is optimized using the learned reward function. Formally, let  $x \sim D_x$  denote prompts and  $y \sim \pi_\theta(\cdot \mid x)$  denote model responses. SFT produces a reference policy  $\pi_{\mathrm{sft}}$, and we set  $\pi_{\mathrm{ref}} := \pi_{\mathrm{sft}}$  unless stated otherwise. A reward model  $r_\phi(x,y)$  is learned from preference comparisons using the Bradley-Terry model \citep{bradley1952rank}. Details are standard and deferred to \Cref{supp:RLHF}.

RLHF then optimizes a KL-regularized objective to prevent reward overoptimization and preserve language quality \citep{Ouyang2022TrainingLM,stiennon2022learningsummarizehumanfeedback,Gao2022ScalingLF,Rafailov2024ScalingLF}:
\begin{align}
  \max_\theta\;\mathbb{E}_{x \sim D_x,\, y \sim \pi_\theta(\cdot\mid x)}\!\big[r_\phi(x,y)\big]\;-\;\beta\,D_{\mathrm{KL}}\!\Big(\pi_\theta(\cdot\mid x)\,\Vert\,\pi_{\mathrm{ref}}(\cdot\mid x)\Big),
\end{align}
equivalently maximizing the KL-regularized reward  $\tilde r(x,y)=r_\phi(x,y)-\beta\log\frac{\pi_\theta(y\mid x)}{\pi_{\mathrm{ref}}(y\mid x)}$. In practice, this objective is optimized with PPO/GRPO/REINFORCE-style methods \citep{williams1992simple,ahmadian2024basicsrevisitingreinforcestyle,shao2024deepseekmathpushinglimitsmathematical,schulman2017proximalpolicyoptimizationalgorithms}.

%%%%%%%%%%%%%%%%%%%%%%%%%%%%%%%%%%%%%%%%%%%%%%%%%%%%%%%%%%%

\subsection{Safe RLHF}
Standard RLHF optimizes a single reward function learned from human preferences, which may be inadequate when attempting to balance competing goals such as \emph{helpfulness} and \emph{harmlessness}. In contrast, Safe RLHF \citep{dai2023safe} separates helpfulness and harmlessness by learning (i) a reward model  $r_\phi(x,y)$, (parameterized by $\phi$ and for input prompt $x$ and model output $y$)  from helpfulness preferences and (ii) a cost model  $c_\psi(x,y)$, (parameterized by $\psi$)  from harmfulness preferences, and then solving the expected-cost constrained RL problem:
\begin{align}
  \max_\theta\;\mathbb{E}\big[r_\phi(x,y)\big] - \beta D_{\mathrm{KL}}(\pi_\theta\Vert \pi_{\mathrm{ref}}) \quad\text{s.t.}\quad \mathbb{E}\big[c_\psi(x,y)\big]\le \tau.
  \label{eq:safe-rlhf-obj}
  % \vspace{-0.3in}f
\end{align}

High-confidence variants like HC-RLHF \citep{chittepu2025reinforcement} replace the expectation constraint with a probabilistic guarantee using the Seldonian framework \citep{thomas2019preventing}.

% Standard RLHF optimizes a single reward function, derived from human preferences, which can be insufficient when trying to balance objectives that may conflict with each other, such as helpfulness and harmlessness. To address this, Safe-RLHF \citep{dai2023safe} imposes an explicit harmfulness constraint, optimizing model responses to improve helpfulness while reducing harmfulness.

% Safe-RLHF collects two preference annotations for each prompt: one indicating which response is more helpful and one indicating which is more harmful. The helpfulness preferences are used to train a reward model, while the harmfulness labels are used to train a cost model. The reward model $r_{\phi}$ and the cost model $c_{\psi}$ are trained using the standard Bradley–Terry objective in Equation~\ref{eq:bt_rm_loss}. Safe-RLHF then solves the following constrained optimization problem to increase the helpfulness of model responses while reducing harmfulness.

% \begin{align}
%   \max_{\theta} \text{\space} &\mathbb{E}_{x \sim \mathcal D_x, y \sim \pi_{\theta}(\cdot \vert x)}[r_{\phi}(x,y)] - \beta \mathbb{D}_\mathrm{KL}\big(\pi_{\theta}(.\vert x) \vert\vert \pi_\mathrm{ref}(. \vert x)\big) \label{eq: obj_safe_rlhf}\\
%   %
%   \text{s.t.} \quad & \mathbb{E}_{x \sim \mathcal D_x, y \sim \pi_{\theta}(\cdot \vert x)}[c_{\psi}(x,y)] \leq \tau. \label{eq: safety constraint}
% \end{align}

% \aj{Can we move this to a related works section?}

% Even though Safe-RLHF reduces the harmfulness of model responses, it does not provide any guarantees on the safety of the resulting model. High Confidence-RLHF (HC-RLHF) \citep{chittepu2025reinforcement} incorporates the Seldonian framework \citep{thomas2019preventing}, modifying the constraint to be a probabilistic upper bound on the expected cost of model responses. It also uses a simpler REINFORCE \citep{ahmadian2024basicsrevisitingreinforcestyle} style algorithm, significantly improving both the helpfulness and harmlessness of model responses compared to Safe-RLHF.

\subsection{First-Order Stochastic Dominance}
\label{sec:fosd}

\aj{cite for FSD} \raj{incorporated Scott's suggestion about describing our approach more accuartely.}
The expected cost constraint in Safe RLFH often proves inadequate in settings where controlling the tail risk is crucial. A more robust method for controlling tail risk is ensuring that the learned policy's cost is stochastically smaller than the cost of the reference policy i.e. less probability must be assigned by the learned policy to high cost outcomes compared to the reference policy. We now describe how we do this formally. Let $F_{X}$ and $Q_{X}$ denote the Cumulative Distribution Function (CDF) and the Quantile Function (Inverse CDF) of a random variable $X$.
For real-valued random variables  $X,Y$, $X$ is said to have first-order stochastic dominance (FSD) on $Y$  (denoted by $X \succeq_{\mathrm{FSD}} Y$)  iff  $F_X(r)\le F_Y(r)$  for all  $r$ \citep{dai2023learning}. Equivalently, we also have $Q_X(q)\ge Q_Y(q)$  for all  $q\in[0,1]$.

%We do this by imposing a first-order stochastic dominance (FSD) constraint. In this section, we decribe the FSD objective formally.

% Since FSD is only a partial order, we use the standard asymmetric violation functional
% \begin{align}
%   \mathcal{L}_{\mathrm{FSD}}(X,Y):=\int_0^1\big(Q_Y(q)-Q_X(q)\big)_+\,dq,
% \end{align}
% which satisfies  $\mathcal{L}_{\mathrm{FSD}}(X,Y)=0$  whenever  $X\succeq_{\mathrm{FSD}} Y$  and relates to Wasserstein-1 via $\mathcal{L}_{\mathrm{FSD}}(X,Y)+\mathcal{L}_{\mathrm{FSD}}(Y,X)=W_1(X,Y)$.
% Moreover, this functional admits an OT characterization with asymmetric convex cost  $c(x,y)=(y-x)_+$  \citep{melnyk2024distributional,Santambrogio2015OptimalTF}, which is what enables our differentiable implementation.

%%%%%%%%%%

% Let $F_{X}$ and $Q_{X}$ denote the Cumulative Distribution Function (CDF) and the Quantile Function (Inverse CDF) of a random variable $X$. We consider the setting when all the random variables are scalar valued, i.e $X, Y \in \mathbb{R}$. A random variable $X$ is said to dominate another random variable $Y$, in First-Order Stochastic Dominance (FSD) sense, $X \succeq_{\text{FSD}} Y$ if

% \begin{equation}
%   F_{X}(r) \leq F_{Y}(r) \quad \forall r \in \mathbb{R}
% \end{equation}

% This can also be expressed in an alternate form, that uses the Quantile function (Inverse CDF) instead of the CDF

% \begin{equation}
%   Q_{X}(q) \geq Q_{Y}(q) \quad \forall q \in [0,1]
% \end{equation}

% One drawback of FSD is that it only imposes a partial ordering over distributions of real valued random variables i.e there can exist two random variables where neither of them dominate the other. This makes it harder to optimize the FSD between two random variables. Hence, we relax the FSD criterion to instead be a loss that measure the extent to which a random variable $X$ \underline{\emph{is dominated by}} another Random Variable $Y$. The FSD loss can be expressed as follows %\aj{since this is a new formulation, should this be in the background section?}

% \begin{equation}
%   \mathcal{L}_\mathrm{FSD}(X,Y) := \int_{0}^{1} \big(Q_{Y}(q) - Q_{X}(q)\big)_{+} dq
%   \label{eq:fsd-loss}
% \end{equation}

% where $(x)_{+} = \max(x,0)$  denotes the ReLU function. This loss measures how much a Random variable $X$ is dominated by $Y$.

Practically, a drawback of FSD is that it only imposes a partial ordering over distributions, so two random variables may be incomparable in terms of FSD. Hence, we relax the definition to an objective that measures the extent to which $X$ falls below $Y$ across quantiles:
\begin{equation}
  \mathcal{L}_{\mathrm{FSD}}(X,Y)
  := \int_{0}^{1} \big(Q_Y(q)-Q_X(q)\big)_+\, dq,
  \label{eq:fsd_loss}
\end{equation}
where $(x)_+=\max(x,0)$ denotes the ReLU function. Note that for two random variables $X$ and $Y$ with different distributions, an FSD objective of $0$ implies that X dominates Y, i.e.  $X \succeq_{\mathrm{FSD}} Y$. More generally,  $\mathcal{L}_{\mathrm{FSD}}(X,Y)$  aggregates the positive quantile gaps  $(Q_Y(q)-Q_X(q))_+$  and therefore quantifies the extent to which  $X$  fails to dominate  $Y$ or $Y$ dominates $X$; it does not by itself certify  $Y\succeq_{\mathrm{FSD}}X$.

In our case, we set $X$ to be the cost distribution $C_{\pi_{\theta}}$ under our learned policy $\Pi_{\theta}$ and $Y$ to be the cost distribution $C_{\pi_{\mathrm{ref}}}$ under the reference policy $\Pi_{\mathrm{ref}}$. Thus, $\mathcal{L}_{\mathrm{FSD}}(C_{\pi_{\theta}},C_{\pi_{\mathrm{ref}}})$ quantifies by how much the cost distribution of the reference policy dominates the cost distribution of the learned policy\footnote{Note that $\mathcal{L}_{\mathrm{FSD}}(X,Y)$ is not a distance because it is not symmetric and can be zero even when $X$ and $Y$ have different distributions. Also $\mathcal{L}_\mathrm{FSD}(X,Y) + \mathcal{L}_\mathrm{FSD}(Y,X) = \mathcal{W}_{1}(X,Y)$ where $\mathcal{W}_{1}(X,Y)$ is the 1-Wassertsein distance }.
% A higher loss implies greater domination by Y on X, but need \underline{\emph{NOT}} imply that Y dominates X.

% \aj{extra: move to footnote}  The following equation relates the FSD objectives $\mathcal{L}_\mathrm{FSD}(X,Y)$ and $\mathcal{L}_\mathrm{FSD}(Y,X)$ to the 1-Wassertsein distance $\mathcal{W}_{1}(X,Y)$:
% \begin{equation}
%   \mathcal{L}_\mathrm{FSD}(X,Y) + \mathcal{L}_\mathrm{FSD}(Y,X) = \mathcal{W}_{1}(X,Y)
% \end{equation}
%In particular, $\mathcal{L}_{\mathrm{FSD}}(X, Y) = 0$ whenever $X \succeq_{\mathrm{FSD}} Y$, even if $X$ and $Y$ are not identically distributed.
%where $\mathcal{W}_{1}(X,Y)$ is the 1-Wasserstein distance.
% between the distributions of X and Y, and can be expressed as \citep{peyre2020computationaloptimaltransport}:
% \begin{equation}
%   \mathcal{W}_{1}(X,Y) = \int_{0}^{1} \big\vert Q_{Y}(q) - Q_{X}(q) \big\vert dq
% \end{equation}

% Using the results from \citet{Santambrogio2015OptimalTF} (Theorem 2.9 and Proposition 2.17) or \citet{melnyk2024distributional} (Theorem 1), the FSD loss can be framed as an optimal transport problem. Specifically,
% \begin{align}
%   \mathcal{L}_\mathrm{FSD}(X,Y) = \mathrm{OT}_c(X,Y)
% \end{align}
% for the asymmetric convex cost function $c(x,y)=(y-x)_+$. We restate the relevant theorem in \Cref{supp:FSD} for completeness.

% As we describe in the next section, FSD is closely related to the optimal transport problem.
%We use \eqref{eq:quantile_particles}--\eqref{eq:empirical_measure} throughout for optimization and gradients.

%With this representation, integrals over quantile levels are approximated by Riemann sums on the fixed grid
%$\{\alpha_i\}_{i=1}^N$.

\subsection{Optimal Transport}\label{subsec:ot}

% Optimal Transport (OT) provides a principled way to compare distributions by optimizing over couplings \citep{peyre2020computationaloptimaltransport}. In our use, OT is primarily a computational tool for building a differentiable surrogate of dominance constraints. We rely on the entropically regularized OT objective \citep{cuturi2013sinkhorn}, which replaces a linear program with a strictly convex problem:
% \begin{align}
%   \mathrm{OT}^{\lambda}_{c}(\mu,\nu) := \min_{P\in \Pi(\mu,\nu)}\;\langle P,C\rangle - \lambda H(P),
% \end{align}
% where  $C_{ij}=c(x_i,y_j)$  and  $H(P)=-\sum_{ij}P_{ij}\log P_{ij}$. The optimizer admits the Sinkhorn scaling form  $P^\star=\mathrm{diag}(u)\,K\,\mathrm{diag}(v)$  with  $K_{ij}=\exp(-C_{ij}/\lambda)$, computable via Sinkhorn iterations \citep{cuturi2013sinkhorn,peyre2020computationaloptimaltransport}. We place full OT/Sinkhorn details in \Cref{supp:ot}, since we only need the fact that the resulting objective is smooth and differentiable w.r.t. its inputs.

%%%%%%%%%%%%%%%%
Optimal transport \citep{peyre2020computationaloptimaltransport} studies the problem of optimally transforming a distribution $\mu$ into another distribution $\nu$ under a cost function
$c(x,y) : \mathcal{X} \times \mathcal{Y} \to \mathbb{R}_+$, which represents the cost of moving a unit mass from $x$ to $y$.

Let $\mu = \sum_{i=1}^N a_i\,\delta_{x_i}$ and $\nu = \sum_{j=1}^M b_j\,\delta_{y_j}$ be two empirical (atomic) distributions on $\mathbb{R}^{d}$, where $x_i, y_j \in \mathbb{R}^d$ are support points (“particles”), $a \in \Delta^N$ and $b \in \Delta^M$ are nonnegative weights summing to 1 ($\Delta^N$ denotes the $N$-dimensional probability simplex), and $\delta_z$ is a Dirac mass at $z$.

A \emph{transport plan} (a.k.a. coupling) from $\mu$ to $\nu$ is a nonnegative matrix $P \in \mathbb{R}_+^{N \times M}$ whose row and column sums match the marginals. Intuitively, $P_{ij}$ represents the amount of probability mass moved from $x_i$ to $y_j$. The set of admissible couplings can be represented as:
\begin{equation}
  \Pi(\mu,\nu) := \Bigl\{ P \in \mathbb{R}_+^{N \times M} \;:\; P \mathbf{1}_M = a, \;\; P^\top \mathbf{1}_N = b \Bigr\}.
  \label{eq:transport-plans}
\end{equation}

Given a ground cost function $c:\mathbb{R}^d \times \mathbb{R}^d \to \mathbb{R}_+$, define the cost matrix $C_{ij} = c(x_i, y_j)$. The \textbf{Kantorovich optimal transport problem} \citep{peyre2020computationaloptimaltransport} between $\mu$ and $\nu$ is
\begin{equation}
  \mathrm{OT}_c(\mu, \nu) := \min_{P \in \Pi(\mu,\nu)} \langle P, C \rangle
  = \min_{P \in \Pi(\mu,\nu)} \sum_{i=1}^N \sum_{j=1}^M P_{ij}\, c(x_i, y_j).
  \label{eq:kantorovich-ot}
\end{equation}

The above objective function is a linear programming problem. Solving this directly is computationally expensive. Hence, to approximate the solution, an entropically regularized OT objective is usually preferred \citep{cuturi2013sinkhorn}, which replaces the linear program with a strictly convex problem:
\begin{align}
  \mathrm{OT}^{\chi}_{c}(\mu,\nu) := \min_{P\in \Pi(\mu,\nu)}\;\langle P,C\rangle - \chi H(P),\label{eq:ot-unreg}
\end{align}
where $H(P)=-\sum_{ij}P_{ij}\log P_{ij}$ and $\chi$ is the regularization parameter. This optimization problem is smooth and has a unique minimizer, which is computable via Sinkhorn iterations \citep{cuturi2013sinkhorn,peyre2020computationaloptimaltransport}. We explain the full details of regularized OT and the Sinkhorn Algorithm in \Cref{supp:ot}.

\yc{Also mention how we are different from Melynk, that they have a direct parametrization of reward with policy, where we only have access to samples and use empirical approximation of distribution i.e our approiach is more general where rewards are generated stochastically and we cannot use the reparameterization trick to express them directly as function of policy parameters}

% \aj{Since the FSD loss is written in terms of integral, should we also write the OT problem in its Lebesgue integral form?}
% \raj{Since we wroking with discrete (empirical) distributions, maybe we can keep this as is and, in section 2.4,define eq 6 as a sum?  }
%\aj{extra} Since we work with (empirical) discrete measures, we use the discrete formulations (see \Cref{eq:fsd_sum_approx}).
%\Cref{eq:kantorovich-ot,eq:ot-unreg}.
%Integrals appear only in the population quantile definitions in \Cref{sec:fosd}, which we approximate by Riemann sums in practice (see \Cref{eq:fsd_sum_approx}).
% \aj{Are the following two equations necessary?}
% A common choice for the cost function is the $p$-norm of the difference between points:
% \begin{equation}
%   c(x_i, y_j) = \| x_i - y_j \|_p^p, \quad p \ge 1,
% \end{equation}
% which leads to the definition of the \emph{$p$-Wasserstein distance} between $\mu$ and $\nu$, denoted as $\mathcal{W}_{p}(\mu, \nu)$.
% \begin{equation}
%   \mathcal{W}_p(\mu, \nu) := \Big(\min_{P \in \Pi(a,b)} \sum_{i=1}^N \sum_{j=1}^M P_{ij}\, \|x_i - y_j\|_p^p.\Big)^{\frac{1}{p}}
%   \label{eq:wasserstein-distance}
% \end{equation}

% \subsubsection{Entropic Regularization and the Sinkhorn Algorithm}

% The optimal transport problem can be solved using linear programming \citep{dantzig2016linear, dvurechensky2018computational}, but this can be computationally expensive when the number of particles is large. Instead, we use entropic regularization, which adds an entropy term to the objective and allows for efficient optimization using the Sinkhorn algorithm \citep{cuturi2013sinkhorn}. Let us define an parameterized empirical distribution $\mu_{\theta}$, where the locations of the particles could be the parameters themselves, and a reference distribution $\nu$. Specifically, we define the entropic regularized OT problem as
% \begin{equation}
%   \mathrm{OT}_{c}^\lambda(\mu_\theta,\nu_{\mathrm{ref}}) \;=\; \min_{P\in\Pi(a,b)}\;\langle P,C\rangle - \lambda H(P), \quad H(P)\;=\;-\sum_{i,j} P_{ij}\log P_{ij},
%   \label{eq:entropic-ot}
% \end{equation}

% for  $\lambda>0$ and a cost matrix $C$. This problem is strictly convex in  $P$  and has a unique minimizer  $P^\star$. More importantly, it can be solved efficiently using the Sinkhorn algorithm \citep{cuturi2013sinkhorn}, which iteratively updates scaling vectors to enforce the marginal constraints. The resulting transport plan  $P^\star$  can then be used to compute the optimal transport objective and its gradient with respect to the policy parameters  $\theta$.

% Define the Gibbs kernel
% \begin{equation}K_{ij} \;=\; e^{-C_{ij}/\lambda}.
%   \label{eq:gibbs-kernel}
% \end{equation}The optimal transport plan can be expressed in closed form as
% \begin{equation}P^\star \;=\; \mathrm{diag}(u) K \mathrm{diag}(v),
%   \label{eq:optimal-plan}
% \end{equation}where  $u \in \mathbb{R}^N$, $v\in\mathbb{R}^M$  are scaling vectors that can be computed using the Sinkhorn iterations
% \begin{equation}u^{(t+1)} \;=\; \frac{a}{Kv^{(t)}}, \qquad v^{(t+1)} \;=\; \frac{b}{K^\top u^{(t+1)}},
%   \label{eq:sinkhorn-iterations}
% \end{equation}starting from some positive initialization (e.g.  $u^{(0)}=v^{(0)}=\mathbf 1$). The iterations are guaranteed to converge to the unique optimal plan  $P^\star$  that satisfies the marginal constraints \citep{peyre2020computationaloptimaltransport, cuturi2013sinkhorn}. Once we have  $P^\star$, we can compute the FSD loss as
% \begin{equation}
%   \mathcal{L}_{\mathrm{FSD}}^\lambda(\mu_\theta,\nu_{\mathrm{ref}}) \;=\; \langle P^\star,C\rangle - \lambda H(P^\star).
%   \label{eq:fsd-loss-sinkhorn}
% \end{equation}

% Because the algorithm consists only of a fixed number of matrix-vector multiplications and divisions, gradients can be propagated through the iterations using automatic differentiation in frameworks like PyTorch \citep{NEURIPS2019_9015} or JAX \citep{jax2018github}, making the Sinkhorn algorithm end to end differentiable.

\paragraph{Relation to FSD.}
We now describe how our FSD objective relates to the optimal transport objective. Using the results from \citet{Santambrogio2015OptimalTF} (Theorem 2.9 and Proposition 2.17) or \citet{melnyk2024distributional} (Theorem 1), for any two distributions $X$ and $Y$, the FSD objective can be framed as an optimal transport problem. Specifically,
\begin{align}\label{eq:otl}
  \mathcal{L}_\mathrm{FSD}(X,Y) = \mathrm{OT}_c(X,Y)
\end{align}
for the asymmetric convex cost function $c(x,y)=(y-x)_+$. We restate the relevant theorem in \Cref{supp:FSD} for completeness. We note that while \cite{melnyk2024distributional} also aims to do alignment via optimal transport, they apply stochastic dominance directly to reward distributions which are derived from a specific parametric form. On the other hand, we only have access to the cost distribution of a policy (independent from the reward distribution) which we access only through sampling. Hence, our framework is more general as we do not assume any specific parametric form. 

% The FSD loss can be framed as an optimal transport problem under an asymmetric convex cost:
% \begin{equation}
%   \mathcal{L}_{\mathrm{FSD}}(X,Y) = \mathrm{OT}_c(\mu_X,\mu_Y),
%   \qquad
%   c(x,y)=(y-x)_+,
%   \label{eq:fsd_ot}
% \end{equation}
% where $\mu_X,\mu_Y$ are the laws of $X,Y$ (see Section C for a formal statement).

% The computation of the FSD loss can be framed as an optimal transport problem, using the results from \citet{Santambrogio2015OptimalTF} (Theorem 2.9 and Proposition 2.17) or \citet{melnyk2024distributional} (Theorem 1). We restate the theorem below.

% \begin{theorem}[from \citet{melnyk2024distributional}]
%   \label{thm:fsd_as_ot}
%   Let $h : \mathbb{R} \to \mathbb{R}_{+}$ be a convex function, and let
%   $X$ and $Y$ be real-valued random variables with probability measures
%   $\mu_X$ and $\mu_Y$, respectively. Denote by
%   $Q_X, Q_Y : [0,1] \to \mathbb{R}$ their (left-continuous) quantile functions.

%   Then,
%   \begin{equation}
%     \int_0^1 h\big(Q_X(q) - Q_Y(q)\big)\, dq
%     =
%     \min_{\gamma \in \Pi(\mu_X,\mu_Y)}
%     \int_{\mathbb{R}^2} h(x-y)\, d\gamma(x,y),
%   \end{equation}
%   where $\Pi(\mu_X,\mu_Y)$ denotes the set of all couplings
%   (joint probability measures) on $\mathbb{R}^2$ with marginals
%   $\mu_X$ and $\mu_Y$. Moreover, if $h$ is strictly convex, the minimizing transport plan $\gamma^\star$ is unique.
% \end{theorem}

% \begin{corollary}[FSD as asymmetric optimal transport]
%   \label{cor:fsd_as_ot}
%   Let $X,Y$ be real-valued random variables with measures
%   $\mu_X$ and $\mu_Y$, and define the FSD loss
%   \[
%     \mathcal{L}_{\mathrm{FSD}}(X,Y)
%     :=
%     \int_0^1 \big(Q_Y(q) - Q_X(q)\big)_+ \, dq.
%   \]
%   Then
%   \[
%     \mathcal{L}_{\mathrm{FSD}}(X,Y)
%     =
%     \min_{\gamma \in \Pi(\mu_X,\mu_Y)}
%     \int_{\mathbb{R}^2} (y-x)_+ \, d\gamma(x,y),
%   \]
%   that is, the FSD loss coincides with the optimal transport cost under the
%   asymmetric convex cost function $c(x,y)=(y-x)_+$.
% \end{corollary}

% Note that the regularized loss  $\mathcal{L}_{\mathrm{FSD}}^\lambda$  is a smooth approximation of the original FSD loss  $\mathcal{L}_{\mathrm{FSD}}$, and as  $\lambda\to0$, we have  $\mathcal{L}_{\mathrm{FSD}}^\lambda \to \mathcal{L}_{\mathrm{FSD}}$. In practice, we can choose a small value of  $\lambda$  to balance approximation accuracy and computational efficiency.

\subsection{Spectral Risk Measures}

Spectral risk measures (SRMs) \citep{acerbi2002spectral} form a broad class of coherent \citep{artzner1999coherent}, law-invariant risk
functionals that aggregate quantile costs using a weighted spectrum over
confidence levels. They provide a flexible way to emphasize parts of the cost distribution (especially the tail) according to application-specific risk tolerance, rather than relying only on the mean. SRMs are widely used in fields like financial risk management and actuarial science to quantify and control exposure to adverse outcomes (\citet{dowd2008spectral,adam2008spectral}). Let $X$ be a real-valued random variable with quantile function
$Q_X : [0,1] \to \mathbb{R}$. A \emph{spectral risk measure} is defined as
\begin{equation}
  \rho_\phi(X)
  =
  \int_0^1 Q_X(q)\,\phi(q)\,dq,
\end{equation}
where $\phi : [0,1] \to \mathbb{R}_+$ is a non-negative, non-decreasing
weighting function satisfying
$\int_0^1 \phi(q)\,dq = 1.$
The function $\phi$ is referred to as the \emph{risk spectrum} and encodes
risk sensitivity across quantiles. Larger weights at higher quantiles emphasize higher
tail costs, while uniform weighting recovers the expectation.
% A canonical example is Conditional Value-at-Risk at level $\alpha$, which corresponds to a step-function spectrum $\phi(q)=\frac{1}{1-\alpha}\mathbf{1}_{[\alpha,1]}(q)$:
% \begin{equation}
%   \mathrm{CVaR}_\alpha(X)
%   =
%   \frac{1}{1-\alpha}
%   \int_\alpha^1 Q_X(q)\,dq.
% \end{equation}
% \aj{Extra}
% SRMs satisfy coherence (monotonicity, subadditivity,
% positive homogeneity, and translation invariance)
% \citep{artzner1999coherent}, and admit the Kusuoka representation,
% which characterizes them as mixtures of Conditional Value-at-Risk (CVaR)
% levels \citep{kusuoka2001law}.

% In Section~\ref{subsec:srm}, we show that how we can use a suitable weighted version of the FSD objective~\eqref{eq:fsd_loss} to represent a broad set of SRMs.
% that places zero weight on quantiles less than $\alpha$ and $\frac{1}{1-\alpha}$ weight everywhere else.
% \begin{equation}
%   \phi(q)
%   =
%   \frac{1}{1-\alpha}\mathbf{1}_{[\alpha,1]}(q).
% \end{equation}

% Thus, SRMs generalize CVaR by allowing structured or smooth weighting over quantiles, providing a flexible mechanism for encoding tail sensitivity in distributional optimization settings.
